\section{System and Method}
\label{sec:method}

\subsection{Overview}
Fig.~\ref{fig:architecture} shows the decision process. An applicant submits an application through the portal.
Eligibility rules are checked first, and the record that will be scored is assembled. A LightGBM model estimates the
probability of repayment difficulty, and a three-band router turns it into Approve, Review or Reject. Every decision
is explained with SHAP and local LIME and stated in plain language. Applicants who are not approved receive recourse
options, and Review cases go to a loan officer. Each component is described below.

\begin{figure*}[t]
\centering
\includegraphics[width=\textwidth]{figures/fig1_architecture.pdf}
\caption{CreditXplain decision process: eligibility rules, risk model, three-band router, explanations for every
decision, recourse for applicants who are not approved, and officer review.}
\label{fig:architecture}
\end{figure*}

\subsection{Eligibility Rules}
\label{sec:eligibility}
Some lending conditions are matters of policy, not of predicted risk. CreditXplain applies them as explicit rules
before the model, so that they are visible and auditable.
\begin{itemize}
\item \emph{Minor applicant.} An applicant under 18 cannot take on a loan alone. The application must include a
guardian (relationship, age, income, income type and years in employment). The record scored by the model takes the
guardian's age, income, income type and employment in place of the applicant's, and every value derived from them is
recomputed. Such an application is never approved automatically, and an officer can approve it only after the
guardian's identity and income documents have been uploaded.
\item \emph{No income.} If the person whose income is assessed has no income, or is recorded as a student or
unemployed, the application is scored but not approved automatically; an officer decides and may ask for a guarantor.
\end{itemize}
A rule only removes automatic approval: an Approve outcome becomes Review, and Review and Reject are unchanged. The
rules that applied, whose details were assessed and the model's own band are stored with the decision and shown in the
reports. The thresholds are configurable policy settings.

\subsection{Risk Model}
The model is LightGBM~\cite{ke2017lightgbm} trained on 300 engineered features built from the application table and
from aggregates of the applicant's credit bureau records and previous applications. The features include combinations
of the three external credit bureau scores, ratios of the loan amount, annuity, goods price and income, employment and
registration histories, and counts of recent credit enquiries. Gender is excluded from the model and kept only for
auditing; age is kept as a feature and audited (Section~\ref{sec:age}). The same feature pipeline is used for
training, evaluation and the running application, so a record scored by the application is identical to one scored
in the experiments.

\subsection{Three-Band Decision Router}
Let $p(x)$ be the predicted probability of repayment difficulty for applicant $x$. The router assigns
\begin{equation}
d(x) =
\begin{cases}
\text{Approve}, & p(x) < \tau_a,\\
\text{Reject}, & p(x) \geq \tau_r,\\
\text{Review}, & \text{otherwise}.
\end{cases}
\label{eq:router}
\end{equation}
Both cutoffs are chosen on the validation set. The approval cutoff $\tau_a$ is the lowest of the 35th, 37th, \ldots,
49th percentiles of the validation scores whose band has a default rate of at most 3\% (the 40th percentile if none
qualifies). The rejection cutoff $\tau_r$ is the 80th percentile, or the 85th if the default rate above the 80th is
below 20\%. An officer may override any automatic decision, but only with a written justification, which is stored in
an audit trail.

\subsection{Explanations}
\emph{SHAP.} For every decision we compute exact SHAP values with TreeExplainer~\cite{lundberg2020tree} on the model's
log-odds. The five largest contributions are shown to the applicant as plain-language reasons, each marked as raising
or lowering risk; every displayed feature has a written plain-language name.

\emph{Local LIME.} Applied to all 300 features with perturbations drawn from the whole training data, LIME produced a
surrogate that described the model's global behaviour rather than the applicant. We restrict it in two ways. Only the
25 features with the highest LightGBM gain are perturbed, while all other features keep the applicant's own values;
and the surrogate is fitted to the model's log-odds, the scale SHAP uses. Perturbations follow LIME's quartile
discretization, with 5{,}000 samples per explanation and a fixed seed, so the same applicant always receives the same
explanation. Officers see both the SHAP contributions and the LIME rules.

\subsection{Recourse Engine}
\label{sec:recourse}
\emph{Levers.} For an applicant who is not approved, the engine looks for changes that would move the application to
the next band: from Reject to Review, or from Review to Approve. It changes only quantities the applicant can act on:
\begin{itemize}
\item the loan amount, scaled by $s_c \in \{1, 0.95, 0.9, 0.85, 0.8, 0.7, 0.6, 0.5\}$, for example through a larger
down payment;
\item documented income, scaled by $s_i \in \{1, 1.1, 1.2, 1.3, 1.5\}$, for example by adding a co-applicant;
\item the yearly repayment, scaled by $s_a \in \{1, 0.9, 0.8, 0.7\}$ through a longer repayment period;
\item waiting $w \in \{0, 3, 6, 12\}$ months before re-applying, while staying in the same job and making no new
credit applications.
\end{itemize}
Waiting is modelled explicitly: employment, registration, identity-document, telephone and car records age by $w$
months, and credit-bureau enquiries move to longer windows or expire. The applicant's own age is held fixed, so no
option depends on growing older. All derived features are recomputed through the same pipeline as the model's input.

\emph{Cost.} An option $s = (s_c, s_i, s_a, w)$ has the cost
\begin{equation}
c(s) = (1 - s_c) + (s_i - 1) + (1 - s_a) + \lambda w,
\label{eq:cost}
\end{equation}
with $\lambda = 0.04$ per month, so that waiting a year costs about as much as halving the loan.

\emph{Safety margin.} An option qualifies only if it reaches the better band with a margin $m$ on the log-odds scale:
\begin{equation}
\operatorname{logit} p(x') \leq \operatorname{logit} \tau - m,
\label{eq:margin}
\end{equation}
where $x'$ is the changed applicant, $\tau$ is the cutoff to be crossed and $m = 0.1$. Without the margin, a
suggestion can sit just below the cutoff and fail after any small change to the model.

\emph{Plausibility.} The changed loan profile (income, loan amount, annuity and goods price, on a log scale and
standardized) must resemble real applicants. We compute its mean distance to the five nearest non-rejected training
applicants and require it to be at most the 95th percentile of the same distance among non-rejected applicants on a
held-back calibration split. This removes combinations such as a large loan with a very small repayment.

\emph{Search.} The engine evaluates the full grid of up to 640 combinations, subject to any levers the applicant has
locked and to the limits the applicant has set on the loan reduction and waiting time. For each combination of levers
it keeps the cheapest qualifying point, then moves that point back towards the applicant's own values in 20 steps and
keeps the cheapest step that still qualifies. Options that ask for at least as much on every lever as another option
are removed, and the three cheapest remaining options are offered, at most one per combination of levers. For each
option, the features that changed are explained with SHAP, so the applicant sees why the change helps.

\subsection{Recourse That Survives Retraining}
\label{sec:robust}
A deployed model is retrained periodically, and an option computed against the current model can fail on the next
one. To anticipate this, we train $K$ bootstrap models: each repeats the training of the final model (same features,
parameters, early stopping and cutoff rule) on a resample, with replacement, of the training data, with its own seed.
An option is called \emph{ensemble-robust} if, in addition to the conditions above, it reaches the better band on all
$K = 5$ of these models, each with its own cutoffs:
\begin{equation}
\left|\{\,k : d_k(x') \text{ is better than } d(x)\,\}\right| \geq K.
\label{eq:agree}
\end{equation}
Requiring agreement makes the search stricter, and some applicants receive no ensemble-robust option. The application
therefore uses a \emph{hybrid}: ensemble-robust options are listed first and labelled as likely to hold if the model is
updated, and standard options fill the remaining places, labelled as possibly changing. The applicant always receives
the same coverage as the standard search, and the most durable advice comes first.

\subsection{Officer Workflow and Reports}
Review cases enter a queue on the officer dashboard. The officer sees the decision, both explanations, the recourse
options, the applicant's limits and any uploaded documents. The officer can approve, reject or request documents; an
override and every document request need a written justification, and every action is recorded with the previous and
new status. Approval produces loan terms priced by risk band. The applicant receives a PDF report with the decision,
the main reasons, the recourse options and the loan terms. The officer can generate a compliance report that adds the
model details, the SHAP and LIME values, the eligibility rules applied, the audit trail and the latest fairness
monitoring results.

The system is implemented as a FastAPI service with an SQLite database and a Next.js web application. The model,
explanation and recourse components are shared by the experiments and the application.
